% Three topologies of a migration between two shard groups of 2f+1 = 5 replicas.
% \topopanel{1} no replication at the recipient, {2} migration layered on replication,
% {3} the integrated chain. \topolegend draws the key shared by the three panels.
%
% Each panel is drawn at 1pt per unit and then scaled by \topozoom so that its fixed
% 178pt bounding box lands on the 0.32\textwidth of a subfigure. The scaling is uniform
% and shared by the legend, so the three panels and the key carry one type size.
\definecolor{ink}{RGB}{34,34,34}
\definecolor{irfill}{gray}{0.86}
\newcommand{\topozoom}{0.927}
\tikzset{topostyle/.style={%
    font=\fontsize{7.5}{9}\fontfamily{phv}\selectfont,
    >={Stealth[length=4pt,width=3.2pt]},
    shard/.style={draw=ink,line width=0.6pt,fill=white,rounded corners=1.5pt,
                  minimum width=26pt,minimum height=14pt,inner sep=0pt},
    leader/.style={shard,fill=ink,text=white,font=\fontsize{7.5}{9}\fontfamily{phv}\bfseries},
    bulk/.style={->,draw=ink,line width=1.8pt,>={Stealth[length=5pt,width=5pt]}},
    abulk/.style={bulk,dash pattern=on 3pt off 2pt},
    notify/.style={->,draw=ink,line width=0.6pt,dash pattern=on 2.2pt off 1.4pt,
                   rounded corners=2pt},
    lbl/.style={inner sep=1pt,text=ink},
    phase/.style={rounded corners=2pt},
    rolelbl/.style={lbl,font=\fontsize{7}{8}\fontfamily{phv}\itshape},
    chunk/.style={draw=ink,line width=0.6pt,fill=irfill,minimum width=18pt,
                  minimum height=10pt,inner sep=0pt,
                  font=\fontsize{6.6}{7.4}\fontfamily{phv}\selectfont},
}}

\newcommand{\topopanel}[1]{%
\scalebox{\topozoom}{%
\begin{tikzpicture}[x=1pt,y=-1pt,topostyle]
  % --- geometry --------------------------------------------------------------
  \def\Dspine{6}   \def\Dx{46}           % donor: log spine, replica column
  \def\Rx{136}     \def\Rspine{174}      % recipient: replica column, log spine
  \def\rowa{22} \def\rowb{48} \def\rowc{74} \def\rowd{100} \def\rowe{126}

  % one bounding box for all three panels, so that they scale alike
  \useasboundingbox (0,12) rectangle (178,166);

  % --- the replicas that make the group durable -----------------------------
  \ifnum#1>1 \fill[phase,irfill] (118,67) rectangle (154,133); \fi

  % --- donor replicas --------------------------------------------------------
  \node[shard]  at (\Dx,\rowa) {A1};
  \node[shard]  at (\Dx,\rowb) {A2};
  \node[shard]  at (\Dx,\rowc) {A3};
  \node[shard]  at (\Dx,\rowd) {A4};
  \node[leader] at (\Dx,\rowe) {A5};
  \node[rolelbl] at (\Dx,141) {leader};

  % --- recipient replicas ----------------------------------------------------
  \node[shard]  at (\Rx,\rowa) {B1};
  \node[shard]  at (\Rx,\rowb) {B2};
  \node[shard]  at (\Rx,\rowc) {B3};
  \node[shard]  at (\Rx,\rowd) {B4};
  \node[leader] at (\Rx,\rowe) {B5};
  \node[rolelbl] at (\Rx,141) {leader};

  % --- replication paths -----------------------------------------------------
  \foreach \y in {\rowa,\rowb,\rowc,\rowd} {
    \draw[notify] (\Dx-13,\rowe) -- (\Dspine,\rowe) -- (\Dspine,\y) -- (\Dx-13,\y);
    \draw[notify] (\Rx+13,\rowe) -- (\Rspine,\rowe) -- (\Rspine,\y) -- (\Rx+13,\y);
  }

  % --- leader-to-leader bulk transfer ---------------------------------------
  \draw[bulk] (\Dx+13,\rowe) -- (\Rx-14,\rowe);
  % The layered panel carries the group as one block. The chain panel carries the
  % same block split in two, so the two spans match end to end.
  \ifnum#1=2
    \node[chunk,minimum width=41pt] at (88,118) {\textit{B}};
  \else
    \node[chunk] at (76.5,118) {\textit{B}/2};
    \node[chunk] at (99.5,118) {\textit{B}/2};
  \fi
  % the brace counts the chunks the group is split into
  \draw[ink,line width=0.5pt,decorate,decoration={brace,amplitude=3pt}]
    (67.5,111) -- (108.5,111);
  \ifnum#1=2
    \node[lbl,anchor=base] at (88,104) {$k = 1$};
  \else
    \node[lbl,anchor=base] at (88,104) {$k = 2$};
  \fi
  \node[lbl,inner sep=1.5pt,
        font=\fontsize{7.5}{9}\fontfamily{phv}\bfseries] at (90,137) {R};

  % --- how the data reaches the followers ------------------------------------
  \ifnum#1=2
    % the leader sends every chunk itself, once per follower
    \foreach \y in {\rowa,\rowb,\rowc,\rowd}
      \draw[bulk,rounded corners=2pt] (\Rx-13,122) -- (114,122) -- (114,\y) -- (\Rx-13,\y);
    \node[lbl,anchor=east,font=\fontsize{7.5}{9}\fontfamily{phv}\bfseries] at (110,61) {R};
  \fi
  % --- what the transfer costs ----------------------------------------------
  \ifnum#1=2
    \node[lbl,anchor=base] at (89,162)
      {$n = 5$:\quad $T_{a} = \tfrac{n+1}{2}\,B/R = 3\,B/R$};
  \else
    \node[lbl,anchor=base] at (89,162)
      {$n = 5$:\quad $T_{b} = \bigl(1 + \tfrac{n-1}{2k}\bigr)B/R = 2\,B/R$};
  \fi

  \ifnum#1=3
    % the chain runs up from the leader, the first two hops reach a majority
    \draw[bulk]  (\Rx,118.5) -- (\Rx,107.5);
    \draw[bulk]  (\Rx,92.5)  -- (\Rx,81.5);
    \node[lbl,anchor=east,font=\fontsize{7.5}{9}\fontfamily{phv}\bfseries] at (120,87) {R};
    \draw[abulk] (\Rx,66.5)  -- (\Rx,55.5);
    \draw[abulk] (\Rx,40.5)  -- (\Rx,29.5);
  \fi
\end{tikzpicture}}%
}

% An empty slot that reserves exactly the space of a panel, for a third figure.
\newcommand{\topoblank}{%
\scalebox{\topozoom}{%
\begin{tikzpicture}[x=1pt,y=-1pt,topostyle]
  \useasboundingbox (0,12) rectangle (178,166);
\end{tikzpicture}}%
}

% The model the two panels above describe, drawn: the time to get one group onto a
% majority of the recipient group, against the number of chunks k, for both ways.
% The formula, not measured data.
%   (a) layered, replication follows the transfer, so splitting does not help:
%       (n+1)/2 B/R, and even a pipelined layered way could not pass (n-1)/2 B/R,
%   (b) chain, every link carries B once over (n+1)/2 pipelined hops:
%       (1 + (n-1)/2k) B/R,
%       so it falls towards B/R.
% The y axis title is turned with \rotatebox rather than with a node rotation, because
% y runs downwards here and a node rotation would come out mirrored.
\newcommand{\topoplot}{%
\scalebox{\topozoom}{%
\begin{tikzpicture}[x=1pt,y=-1pt,topostyle,
    axis/.style={draw=ink,line width=0.7pt},
    ticklbl/.style={inner sep=1.5pt,text=ink,
                    font=\fontsize{6.6}{7.4}\fontfamily{phv}\selectfont},
    layered/.style={draw=ink,line width=1.6pt,dash pattern=on 4pt off 2.5pt},
    chain/.style={draw=ink,line width=1.6pt,line join=round},
    grid/.style={draw=ink!45,line width=0.5pt,dash pattern=on 1.2pt off 1.6pt}]
  \useasboundingbox (0,12) rectangle (178,166);

  \node[ticklbl,inner sep=0pt] at (14,82.0)
    {\rotatebox{90}{migration transfer duration ($B/R$)}};

  % the floors each way is bound by, doubling as the gridlines
  \draw[grid] (32.0,70.5) -- (172.0,70.5);
  \node[ticklbl,anchor=north east] at (171.0,71.5) {$\tfrac{n-1}{2}B/R$};
  \draw[grid] (32.0,103.3) -- (172.0,103.3);
  \node[ticklbl,anchor=north east] at (171.0,104.3) {$B/R$};

  % axes
  \draw[axis] (32.0,136.0) -- (172.0,136.0);
  \draw[axis] (32.0,136.0) -- (32.0,28.0);
  \draw[axis] (32.0,136.0) -- (32.0,139.0);
  \node[ticklbl,anchor=north] at (32.0,140.0) {1};
  \draw[axis] (50.9,136.0) -- (50.9,139.0);
  \node[ticklbl,anchor=north] at (50.9,140.0) {2};
  \draw[axis] (88.6,136.0) -- (88.6,139.0);
  \node[ticklbl,anchor=north] at (88.6,140.0) {4};
  \draw[axis] (126.3,136.0) -- (126.3,139.0);
  \node[ticklbl,anchor=north] at (126.3,140.0) {6};
  \draw[axis] (164.0,136.0) -- (164.0,139.0);
  \node[ticklbl,anchor=north] at (164.0,140.0) {8};
  \draw[axis] (32.0,103.3) -- (29.0,103.3);
  \node[ticklbl,anchor=east] at (27.0,103.3) {1};
  \draw[axis] (32.0,70.5) -- (29.0,70.5);
  \node[ticklbl,anchor=east] at (27.0,70.5) {2};
  \draw[axis] (32.0,37.8) -- (29.0,37.8);
  \node[ticklbl,anchor=east] at (27.0,37.8) {3};
  \node[ticklbl,anchor=north] at (98.0,152.0) {chunks $k$};

  \node[ticklbl,anchor=north west] at (38,112) {$n = 5$};

  % (a) layered: k does not help, (1+f) B/R
  \draw[layered] (32.0,37.8) -- (164.0,37.8);
  \node[ticklbl,anchor=south east] at (164.0,36.3) {(a) layered};

  % (b) chain: (1 + f/k) B/R
  \draw[chain] (32.0,37.8) -- (50.9,70.5) -- (69.7,81.5) -- (88.6,86.9) -- (107.4,90.2) -- (126.3,92.4) -- (145.1,93.9) -- (164.0,95.1);
  \fill[ink] (32.0,37.8) circle (1.4pt);
  \fill[ink] (50.9,70.5) circle (1.4pt);
  \fill[ink] (69.7,81.5) circle (1.4pt);
  \fill[ink] (88.6,86.9) circle (1.4pt);
  \fill[ink] (107.4,90.2) circle (1.4pt);
  \fill[ink] (126.3,92.4) circle (1.4pt);
  \fill[ink] (145.1,93.9) circle (1.4pt);
  \fill[ink] (164.0,95.1) circle (1.4pt);
  \node[ticklbl,anchor=south] at (107.4,86.2) {(b) chain};
\end{tikzpicture}}%
}

\newcommand{\topolegend}{%
\scalebox{\topozoom}{%
\begin{tikzpicture}[x=1pt,y=-1pt,topostyle]
  \def\key##1{\tikz[baseline=-2.7pt]\draw[##1] (0pt,0pt) -- (18pt,0pt);\hspace{4pt}}
  \node[inner sep=0pt,anchor=base] at (0,0) {%
    \key{notify}Replicated log entry\hspace{18pt}%
    \key{bulk}Synchronous RDMA write\hspace{18pt}%
    \key{abulk}Asynchronous RDMA write};
\end{tikzpicture}}%
}
